\renewcommand{\thesection}{\arabic{section}}\renewcommand{\theHsection}{chapter.\arabic{section}}\setcounter{section}{12}\section{反常积分和含参变量的积分}
\noindent\begin{minipage}{\textwidth}\ProblemHeading{例题}{13.1}
证明：对任意实数 $\alpha$，反常积分 $\displaystyle\int_1^{+\infty}x^\alpha e^{-x}\,dx$ 收敛。
\end{minipage}\par\vfill
\noindent\begin{minipage}{\textwidth}\ProblemHeading{例题}{13.2}
证明反常积分 $\displaystyle\int_1^{+\infty}\frac{x^2}{\sqrt{x+1}(x^4+x+1)}\,dx$ 收敛。
\end{minipage}\par\vfill
\noindent\begin{minipage}{\textwidth}\ProblemHeading{例题}{13.3}
设 $g$ 在 $[a,+\infty)$ 上单调且 $\lim_{x\to+\infty}g(x)=0$。证明下列积分收敛：
\SubQ{1}{$\displaystyle\int_a^{+\infty}g(x)\sin x\,dx$。}
\SubQ{2}{$\displaystyle\int_a^{+\infty}g(x)\cos x\,dx$。}
\end{minipage}\par\vfill
\noindent\begin{minipage}{\textwidth}\ProblemHeading{例题}{13.4}
设 $a>0$，$0<p\le1$。证明下列反常积分条件收敛：
\SubQ{1}{$\displaystyle\int_a^{+\infty}\frac{\sin x}{x^p}\,dx$。}
\SubQ{2}{$\displaystyle\int_a^{+\infty}\frac{\cos x}{x^p}\,dx$。}
\end{minipage}\par\vfill
\newpage
\noindent\begin{minipage}{\textwidth}\ProblemHeading{例题}{13.5}
设 $0<k<1$。研究椭圆积分 $\displaystyle\int_0^1\frac{dx}{\sqrt{(1-x^2)(1-k^2x^2)}}$ 的敛散性。
\end{minipage}\par\vfill
\noindent\begin{minipage}{\textwidth}\ProblemHeading{例题}{13.6}
研究反常积分 $\displaystyle\int_0^1\frac{\ln x}{1-x^2}\,dx$ 的敛散性。
\end{minipage}\par\vfill
\noindent\begin{minipage}{\textwidth}\ProblemHeading{例题}{13.7}
研究反常积分 $\displaystyle\int_0^{+\infty}\frac{x^\alpha\arctan x}{2+x^\beta}\,dx$ 的敛散性，其中 $\alpha\in\mathbb R$，$\beta\ge0$。
\end{minipage}\par\vfill
\noindent\begin{minipage}{\textwidth}\ProblemHeading{习题}{13.8}
判断下列反常积分的敛散性（含参数时按参数讨论）：
\SubQ{1}{$\displaystyle\int_0^{+\infty}\frac{\ln(x^2+1)}x\,dx$。}
\SubQ{2}{$\displaystyle\int_0^{+\infty}\sqrt{x}e^{-x}\,dx$。}
\SubQ{3}{$\displaystyle\int_0^{+\infty}\frac{x\arctan x}{\sqrt[3]{1+x^4}}\,dx$。}
\end{minipage}\par\vfill
\newpage
\noindent\begin{minipage}{\textwidth}\ProblemHeading{习题}{13.9}
判断下列反常积分的敛散性（含参数时按参数讨论）：
\SubQ{1}{$\displaystyle\int_{e^2}^{+\infty}\frac{dx}{x\ln\ln x}$。}
\SubQ{2}{$\displaystyle\int_0^1\frac{\ln x}{\sqrt{1-x^2}}\,dx$。}
\SubQ{3}{$\displaystyle\int_0^1\frac{x^2}{\sqrt[3]{(1-x^2)^5}}\,dx$。}
\end{minipage}\par\vfill
\noindent\begin{minipage}{\textwidth}\ProblemHeading{习题}{13.10}
判断下列反常积分的敛散性（含参数时按参数讨论）：
\SubQ{1}{$\displaystyle\int_0^1\frac{dx}{e^{\sqrt{x}}-1}$。}
\SubQ{2}{$\displaystyle\int_0^1\frac{\sqrt{x}}{e^{\sin x}-1}\,dx$。}
\SubQ{3}{$\displaystyle\int_0^1\frac{dx}{e^x-\cos x}$。}
\end{minipage}\par\vfill
\noindent\begin{minipage}{\textwidth}\ProblemHeading{习题}{13.11}
判断下列反常积分的敛散性（含参数时按参数讨论）：
\SubQ{1}{$\displaystyle\int_0^{\pi/2}\frac{\ln\sin x}{\sqrt{x}}\,dx$。}
\SubQ{2}{$\displaystyle\int_0^{\pi/2}\frac{dx}{\sqrt{\sin x}\cos x}$。}
\SubQ{3}{$\displaystyle\int_0^{+\infty}\frac{\arctan x}{x^\mu}\,dx$，$\mu\in\mathbb R$。}
\end{minipage}\par\vfill
\noindent\begin{minipage}{\textwidth}\ProblemHeading{习题}{13.12}
研究下列积分的敛散性，并区分绝对收敛与条件收敛：
\SubQ{1}{$\displaystyle\int_1^{+\infty}\frac{\cos(1-2x)}{\sqrt[3]x\sqrt[3]{x^2+1}}\,dx$。}
\SubQ{2}{$\displaystyle\int_0^{+\infty}\frac{\sin x}{\sqrt[3]{x^2+x+1}}\,dx$。}
\SubQ{3}{$\displaystyle\int_2^{+\infty}\frac{\sin x}{x\ln x}\,dx$。}
\end{minipage}\par\vfill
\newpage
\noindent\begin{minipage}{\textwidth}\ProblemHeading{习题}{13.13}
研究下列积分的敛散性，并区分绝对收敛与条件收敛：
\SubQ{1}{$\displaystyle\int_0^{+\infty}\frac{\ln(1+x)}{x^2(1+x^p)}\,dx$，$p>0$。}
\SubQ{2}{$\displaystyle\int_0^{+\infty}\frac{\sin x}{x(1+\sqrt{x})}\,dx$。}
\SubQ{3}{$\displaystyle\int_0^1\frac{\sin(1/x)}{x^p}\,dx$，$p\in\mathbb R$。}
\end{minipage}\par\vfill
\noindent\begin{minipage}{\textwidth}\ProblemHeading{习题}{13.14}
设 $f$ 在 $[a,+\infty)$ 上单调、连续，且 $\displaystyle\int_a^{+\infty}f(x)\,dx$ 收敛。证明 $\displaystyle\lim_{x\to+\infty}f(x)=0$。
\end{minipage}\par\vfill
\noindent\begin{minipage}{\textwidth}\ProblemHeading{习题}{13.15}
设 $f,g$ 在 $[0,+\infty)$ 上非负，$\displaystyle\int_0^{+\infty}g(x)\,dx$ 收敛，且当 $0<x<y$ 时有 $\displaystyle f(y)\le f(x)+\int_x^y g(t)\,dt$。证明 $\displaystyle\lim_{x\to+\infty}f(x)$ 存在。
\end{minipage}\par\vfill
\noindent\begin{minipage}{\textwidth}\ProblemHeading{习题}{13.16}
设 $f,g$ 在 $[0,+\infty)$ 上非负，并在有限区间上可积；$g$ 单调递减趋于 $0$，且 $\displaystyle\int_0^{+\infty}f(x)g(x)\,dx$ 收敛。证明 $\displaystyle\lim_{x\to+\infty}g(x)\int_0^x f(t)\,dt=0$。
\end{minipage}\par\vfill
\newpage
\noindent\begin{minipage}{\textwidth}\ProblemHeading{习题}{13.17}
设 $\displaystyle\int_0^{+\infty}f(x)\,dx$ 绝对收敛，$g$ 在有限区间上可积且 $\displaystyle\lim_{x\to+\infty}g(x)=0$。证明 $\displaystyle\lim_{x\to+\infty}\int_0^x f(t)g(x-t)\,dt=0$。
\end{minipage}\par\vfill
\noindent\begin{minipage}{\textwidth}\ProblemHeading{例题}{13.18}
计算反常二重积分 $\displaystyle\iint_{\mathbb R^2}e^{-x^2-y^2}\,dx\,dy$。
\end{minipage}\par\vfill
\noindent\begin{minipage}{\textwidth}\ProblemHeading{例题}{13.19（改写）}
设实数项级数 $\sum_{n=1}^{\infty}a_n$ 条件收敛，并令 $f(x)=a_n$（$n-1\le x<n$，$n=1,2,\ldots$）。比较下面两种定义下反常积分的敛散性，并说明理由。
\SubQ{1}{通常的单变量反常积分 $\displaystyle\lim_{A\to+\infty}\int_0^A f(x)\,dx$。}
\SubQ{2}{要求对任意有界 Jordan 可测集的递增列 $D_m\subset[0,+\infty)$，只要 $\bigcup_{m=1}^{\infty}D_m=[0,+\infty)$，积分 $\int_{D_m}f(x)\,dx$ 都趋向同一个有限极限。}
\end{minipage}\par\vfill
\noindent\begin{minipage}{\textwidth}\ProblemHeading{习题}{13.20}
计算下列反常二重积分：
\SubQ{1}{$\displaystyle\iint_D\frac{dx\,dy}{\sqrt{x^2+y^2}}$，其中 $D=\{(x,y):x^2+y^2<1\}$。}
\SubQ{2}{$\displaystyle\iint_D\frac{dx\,dy}{(1+x+y)^\alpha}$，其中 $D=\{x>0,y>0\}$，$\alpha>2$。}
\SubQ{3}{$\displaystyle\iint_D\max\{x,y\}e^{-(x^2+y^2)}\,dx\,dy$，其中 $D=\{x>0,y>0\}$。}
\end{minipage}\par\vfill
\newpage
\noindent\begin{minipage}{\textwidth}\ProblemHeading{习题}{13.21}
已知 Fresnel 积分 $\displaystyle\int_{-\infty}^{+\infty}\sin x^2\,dx=\int_{-\infty}^{+\infty}\cos x^2\,dx=\sqrt{\pi/2}$。验证 $\displaystyle\int_{-\infty}^{+\infty}dy\int_{-\infty}^{+\infty}\sin(x^2+y^2)\,dx=\int_{-\infty}^{+\infty}dx\int_{-\infty}^{+\infty}\sin(x^2+y^2)\,dy=\pi$。再证明 $\sin(x^2+y^2)$ 在 $\mathbb R^2$ 上的反常二重积分不收敛：这里收敛要求沿任意有界 Jordan 可测区域的穷尽递增列所得积分极限存在且相同。
\end{minipage}\par\vfill
\noindent\begin{minipage}{\textwidth}\ProblemHeading{例题}{13.22}
计算积分 $\displaystyle\int_0^1\frac{\ln(1+x)}{1+x^2}\,dx$。
\end{minipage}\par\vfill
\noindent\begin{minipage}{\textwidth}\ProblemHeading{例题}{13.23}
设 $\displaystyle I(u)=\int_u^{u^2}\frac{\sin(ux)}x\,dx$，在 $x=0$ 处按连续延拓理解被积函数。求 $I'(u)$。
\end{minipage}\par\vfill
\noindent\begin{minipage}{\textwidth}\ProblemHeading{习题}{13.24}
试用两种方法计算下列极限：
\SubQ{1}{$\displaystyle\lim_{\alpha\to0}\int_{-1}^1\sqrt{x^2+\alpha^2}\,dx$。}
\SubQ{2}{$\displaystyle\lim_{\alpha\to0}\int_\alpha^{1+\alpha}\frac{dx}{1+x^2+\alpha^2}$。}
\end{minipage}\par\vfill
\newpage
\noindent\begin{minipage}{\textwidth}\ProblemHeading{习题}{13.25}
求下列函数的导数 $F'(\alpha)$（可去间断点均作连续延拓）：
\SubQ{1}{$\displaystyle F(\alpha)=\int_{\sin\alpha}^{\cos\alpha}e^{\alpha\sqrt{1-x^2}}\,dx$。}
\SubQ{2}{$\displaystyle F(\alpha)=\int_{a+\alpha}^{b+\alpha}\frac{\sin(\alpha x)}x\,dx$，其中 $a,b$ 为常数。}
\SubQ{3}{$\displaystyle F(\alpha)=\int_0^\alpha\frac{\ln(1+\alpha x)}x\,dx$。}
\end{minipage}\par\vfill
\noindent\begin{minipage}{\textwidth}\ProblemHeading{习题}{13.26}
求下列函数的导数 $F'(\alpha)$（可去间断点均作连续延拓）：
\SubQ{1}{$\displaystyle F(\alpha)=\int_0^\alpha f(x+\alpha,x-\alpha)\,dx$，其中 $f(u,v)$ 有连续偏导数。}
\end{minipage}\par\vfill
\noindent\begin{minipage}{\textwidth}\ProblemHeading{习题}{13.27}
设 $f$ 在 $[a,b]$ 上连续，$c\in[a,b]$，$k\ne0$ 为常数。证明 $\displaystyle y(x)=\frac1k\int_c^x f(t)\sin k(x-t)\,dt$ 在 $[a,b]$ 上满足常微分方程 $y''+k^2y=f(x)$。\textbf{（校订：补明 $k\ne0$。）}
\end{minipage}\par\vfill
\noindent\begin{minipage}{\textwidth}\ProblemHeading{习题}{13.28}
应用对参数进行微分或积分的方法计算：
\SubQ{1}{$\displaystyle\int_0^{\pi/2}\ln(a^2\sin^2x+b^2\cos^2x)\,dx$，$a,b>0$。}
\SubQ{2}{$\displaystyle\int_0^\pi\ln(1-2a\cos x+a^2)\,dx$，$0\le a<1$。}
\SubQ{3}{$\displaystyle\int_0^{\pi/2}\frac{\arctan(a\tan x)}{\tan x}\,dx$，$a\ge0$。}
\end{minipage}\par\vfill
\newpage
\noindent\begin{minipage}{\textwidth}\ProblemHeading{习题}{13.29}
应用对参数进行微分或积分的方法计算：
\SubQ{1}{$\displaystyle\int_0^{\pi/2}\ln\frac{1+a\cos x}{1-a\cos x}\frac{dx}{\cos x}$，$0\le a<1$。}
\end{minipage}\par\vfill
\noindent\begin{minipage}{\textwidth}\ProblemHeading{习题}{13.30}
设 $f(x,u)$ 在 $[a,+\infty)\times[\alpha,\beta]$ 上连续。证明以下含参反常积分判别法：
\SubQ{1}{Cauchy 准则：$\int_a^{+\infty}f(x,u)\,dx$ 关于 $u\in[\alpha,\beta]$ 一致收敛，当且仅当任给 $\varepsilon>0$，存在 $X$，使对一切 $A',A''>X$ 及 $u\in[\alpha,\beta]$ 都有 $\displaystyle\left|\int_{A'}^{A''}f(x,u)\,dx\right|<\varepsilon$。}
\SubQ{2}{Weierstrass 判别法：若有在 $[a,+\infty)$ 上可积的函数 $p$，使对一切充分大的 $x$ 及所有 $u\in[\alpha,\beta]$ 有 $|f(x,u)|\le p(x)$，则 $\displaystyle\int_a^{+\infty}f(x,u)\,dx$ 在 $[\alpha,\beta]$ 上一致收敛。}
\end{minipage}\par\vfill
\noindent\begin{minipage}{\textwidth}\ProblemHeading{例题}{13.31}
设 $\alpha>0$。证明：
\SubQ{1}{$\displaystyle\int_0^{+\infty}e^{-(\alpha+u^2)x}\sin x\,dx$ 关于 $u\in[0,+\infty)$ 一致收敛。}
\SubQ{2}{$\displaystyle\int_0^{+\infty}e^{-(\alpha+u^2)x}\sin x\,du$ 关于 $x\in[0,+\infty)$ 一致收敛。}
\end{minipage}\par\vfill
\noindent\begin{minipage}{\textwidth}\ProblemHeading{例题}{13.32}
证明积分 $\displaystyle\int_a^{+\infty}ue^{-ux}\,dx$ 关于 $u\in[0,+\infty)$ 不一致收敛，其中 $a$ 为实常数。
\end{minipage}\par\vfill
\newpage
\noindent\begin{minipage}{\textwidth}\ProblemHeading{例题}{13.33}
设 $\alpha>0$，$\beta_0>0$。证明积分 $\displaystyle\int_0^{+\infty}\frac{x\sin(\beta x)}{\alpha^2+x^2}\,dx$ 关于 $\beta\in[\beta_0,+\infty)$ 一致收敛。
\end{minipage}\par\vfill
\noindent\begin{minipage}{\textwidth}\ProblemHeading{例题}{13.34}
设 $0<p\le1$。证明积分 $\displaystyle\int_0^{+\infty}e^{-ux}\frac{\sin x}{x^p}\,dx$ 关于 $u\in[0,+\infty)$ 一致收敛。
\end{minipage}\par\vfill
\noindent\begin{minipage}{\textwidth}\ProblemHeading{例题}{13.35}
设 $0<a<b$。计算反常积分 $\displaystyle\int_0^{+\infty}\frac{e^{-ax}-e^{-bx}}x\,dx$。
\end{minipage}\par\vfill
\noindent\begin{minipage}{\textwidth}\ProblemHeading{例题}{13.36}
计算含参反常积分 $\displaystyle\int_0^{+\infty}e^{-x^2}\cos(2\beta x)\,dx$，其中 $\beta\in\mathbb R$。
\end{minipage}\par\vfill
\newpage
\noindent\begin{minipage}{\textwidth}\ProblemHeading{例题}{13.37（改写）}
利用含参反常积分的性质，计算 Dirichlet 积分 $\displaystyle\int_0^{+\infty}\frac{\sin x}{x}\,dx$。
\end{minipage}\par\vfill
\noindent\begin{minipage}{\textwidth}\ProblemHeading{例题}{13.38（改写）}
利用含参反常积分的性质，计算下列 Laplace 积分：
\SubQ{1}{$\displaystyle\int_0^{+\infty}\frac{\cos(\beta x)}{\alpha^2+x^2}\,dx$，其中 $\alpha>0$，$\beta\ge0$。}
\SubQ{2}{$\displaystyle\int_0^{+\infty}\frac{x\sin(\beta x)}{\alpha^2+x^2}\,dx$，其中 $\alpha>0$，$\beta>0$。}
\end{minipage}\par\vfill
\noindent\begin{minipage}{\textwidth}\ProblemHeading{例题}{13.39（改写）}
计算下列 Fresnel 积分：
\SubQ{1}{$\displaystyle\int_0^{+\infty}\sin(x^2)\,dx$。}
\SubQ{2}{$\displaystyle\int_0^{+\infty}\cos(x^2)\,dx$。}
\end{minipage}\par\vfill
\noindent\begin{minipage}{\textwidth}\ProblemHeading{习题}{13.40}
求下列含参反常积分收敛时参数的取值范围：
\SubQ{1}{$\displaystyle\int_0^{+\infty}x^u\,dx$。}
\SubQ{2}{$\displaystyle\int_1^{+\infty}x^u\frac{x+\sin x}{x-\sin x}\,dx$。}
\SubQ{3}{$\displaystyle\int_2^{+\infty}\frac{dx}{x^u\ln x}$。}
\end{minipage}\par\vfill
\newpage
\noindent\begin{minipage}{\textwidth}\ProblemHeading{习题}{13.41}
求下列含参反常积分收敛时参数的取值范围：
\SubQ{1}{$\displaystyle\int_0^\pi\frac{dx}{\sin^u x}$。}
\SubQ{2}{$\displaystyle\int_0^{+\infty}\frac{\sin^2x}{x^\alpha(1+x)}\,dx$。}
\SubQ{3}{$\displaystyle\int_0^{+\infty}\frac{\ln(1+x^2)}{x^\alpha}\,dx$。}
\end{minipage}\par\vfill
\noindent\begin{minipage}{\textwidth}\ProblemHeading{习题}{13.42}
研究下列积分在指定参数范围内的一致收敛性：
\SubQ{1}{$\displaystyle\int_{-\infty}^{+\infty}\frac{\cos(ux)}{1+x^2}\,dx$，$u\in\mathbb R$。}
\SubQ{2}{$\displaystyle\int_0^{+\infty}e^{-\alpha x}\sin(\beta x)\,dx$，$\alpha\ge\alpha_0>0$；$\beta$ 为固定实数。}
\SubQ{3}{$\displaystyle\int_0^{+\infty}e^{-\alpha x}\sin(\beta x)\,dx$，$\alpha>0$；$\beta$ 为固定实数（区分 $\beta=0$ 与 $\beta\ne0$）。}
\end{minipage}\par\vfill
\noindent\begin{minipage}{\textwidth}\ProblemHeading{习题}{13.43}
研究下列积分在指定参数范围内的一致收敛性：
\SubQ{1}{$\displaystyle\int_0^{+\infty}\sqrt\alpha\,e^{-\alpha x^2}\,dx$，$\alpha\ge0$。}
\SubQ{2}{$\displaystyle\int_1^{+\infty}\frac{\ln(1+x^2)}{x^\alpha}\,dx$，$\alpha>1$。}
\SubQ{3}{$\displaystyle\int_1^{+\infty}e^{-\alpha x}\frac{\cos x}{x^p}\,dx$，$\alpha\ge0$，其中 $p>0$ 为常数。}
\end{minipage}\par\vfill
\noindent\begin{minipage}{\textwidth}\ProblemHeading{习题}{13.44}
研究下列积分在指定参数范围内的一致收敛性：
\SubQ{1}{$\displaystyle\int_0^{+\infty}\frac{\sin(x^2)}{1+x^p}\,dx$，$0\le p<+\infty$。}
\end{minipage}\par\vfill
\newpage
\noindent\begin{minipage}{\textwidth}\ProblemHeading{习题}{13.45}
设 $f(x,u)$ 在 $a\le x<+\infty$、$\alpha\le u\le\beta$ 上连续，且对每个 $u\in[\alpha,\beta)$，积分 $\displaystyle\int_a^{+\infty}f(x,u)\,dx$ 收敛，而 $\displaystyle\int_a^{+\infty}f(x,\beta)\,dx$ 发散。证明该含参积分在 $[\alpha,\beta)$ 上不一致收敛。
\end{minipage}\par\vfill
\noindent\begin{minipage}{\textwidth}\ProblemHeading{习题}{13.46}
设 $a$ 为实常数。证明 $\displaystyle\varphi(u)=\int_a^{+\infty}ue^{-ux}\,dx$ 在 $u>0$ 上连续，并说明为什么这与该积分在 $u\ge0$ 上不一致收敛并不矛盾。
\end{minipage}\par\vfill
\noindent\begin{minipage}{\textwidth}\ProblemHeading{习题}{13.47}
验证下列两个累次积分不相等，并说明理由：\[\int_0^1du\int_0^{+\infty}(2-xu)xu e^{-xu}\,dx\ne\int_0^{+\infty}dx\int_0^1(2-xu)xu e^{-xu}\,du.\]
\end{minipage}\par\vfill
\noindent\begin{minipage}{\textwidth}\ProblemHeading{习题}{13.48}
证明 $\displaystyle F(\alpha)=\int_0^{+\infty}\frac{\cos x}{1+(x+\alpha)^2}\,dx$ 在 $0\le\alpha<+\infty$ 上连续且可微（端点取右导数）。
\end{minipage}\par\vfill
\newpage
\noindent\begin{minipage}{\textwidth}\ProblemHeading{习题}{13.49}
计算下列反常积分：
\SubQ{1}{$\displaystyle\int_0^1\frac{x^\beta-x^\alpha}{\ln x}\,dx$，$\alpha,\beta>-1$。}
\SubQ{2}{$\displaystyle\int_0^{+\infty}\frac{1-e^{-ax}}{xe^x}\,dx$，$a>-1$。}
\SubQ{3}{$\displaystyle\int_0^{+\infty}\frac{1-e^{-ax^2}}{x^2}\,dx$，$a>0$。}
\end{minipage}\par\vfill
\noindent\begin{minipage}{\textwidth}\ProblemHeading{习题}{13.50}
计算下列反常积分：
\SubQ{1}{$\displaystyle\int_0^{+\infty}\frac{e^{-\alpha x^2}-e^{-\beta x^2}}x\,dx$，$\alpha,\beta>0$。}
\SubQ{2}{$\displaystyle\int_0^{+\infty}\frac{\arctan(ax)}{x(1+x^2)}\,dx$，$a>0$。}
\SubQ{3}{$\displaystyle\int_0^{+\infty}\left[e^{-(a/x)^2}-e^{-(b/x)^2}\right]\,dx$，$0<a<b$。}
\end{minipage}\par\vfill
\noindent\begin{minipage}{\textwidth}\ProblemHeading{习题}{13.51}
利用 $\displaystyle\int_0^{+\infty}e^{-x^2}\,dx=\frac{\sqrt\pi}2$ 及 $\displaystyle\int_0^{+\infty}\frac{\sin x}x\,dx=\frac\pi2$，计算：
\SubQ{1}{$\displaystyle\int_{-\infty}^{+\infty}\frac{x}{\sigma\sqrt{2\pi}}e^{-\frac12((x-a)/\sigma)^2}\,dx$，$\sigma>0$，$a\in\mathbb R$。}
\SubQ{2}{$\displaystyle\int_{-\infty}^{+\infty}\frac{(x-a)^2}{\sigma\sqrt{2\pi}}e^{-\frac12((x-a)/\sigma)^2}\,dx$，$\sigma>0$，$a\in\mathbb R$。}
\SubQ{3}{$\displaystyle\int_0^{+\infty}\frac{\sin(ax)\cos(bx)}x\,dx$，$a,b>0$。}
\end{minipage}\par\vfill
\noindent\begin{minipage}{\textwidth}\ProblemHeading{习题}{13.52}
利用 $\displaystyle\int_0^{+\infty}e^{-x^2}\,dx=\frac{\sqrt\pi}2$ 及 $\displaystyle\int_0^{+\infty}\frac{\sin x}x\,dx=\frac\pi2$，计算：
\SubQ{1}{$\displaystyle\int_0^{+\infty}\frac{\sin^2x}{x^2}\,dx$。}
\SubQ{2}{$\displaystyle\int_0^{+\infty}x^{2n}e^{-x^2}\,dx$，$n$ 为正整数。}
\SubQ{3}{$\displaystyle\int_0^{+\infty}\frac{\sin^4x}{x^2}\,dx$。}
\end{minipage}\par\vfill
\newpage
\noindent\begin{minipage}{\textwidth}\ProblemHeading{例题}{13.53}
求积分 $\displaystyle\int_0^{\pi/2}\sin^n x\cos^m x\,dx$ 的值，其中 $n,m$ 均为非负整数。
\end{minipage}\par\vfill
\noindent\begin{minipage}{\textwidth}\ProblemHeading{习题}{13.54}
设 $\displaystyle\Gamma(s)=\int_0^{+\infty}t^{s-1}e^{-t}\,dt$（$s>0$）。证明：
\SubQ{1}{$\displaystyle\Gamma(s)=2\int_0^{+\infty}x^{2s-1}e^{-x^2}\,dx$，$s>0$。}
\SubQ{2}{$\displaystyle\Gamma(s)=a^s\int_0^{+\infty}x^{s-1}e^{-ax}\,dx$，$s,a>0$。}
\end{minipage}\par\vfill
\noindent\begin{minipage}{\textwidth}\ProblemHeading{习题}{13.55}
设 $\displaystyle B(p,q)=\int_0^1t^{p-1}(1-t)^{q-1}\,dt$。证明当 $p,q>0$ 时，$\displaystyle B(p,q)=2\int_0^{\pi/2}\sin^{2p-1}t\cos^{2q-1}t\,dt$。
\end{minipage}\par\vfill
\noindent\begin{minipage}{\textwidth}\ProblemHeading{习题}{13.56}
利用 Euler 积分计算：
\SubQ{1}{$\displaystyle\int_0^1\sqrt{x-x^2}\,dx$。}
\SubQ{2}{$\displaystyle\int_0^a x^2\sqrt{a^2-x^2}\,dx$，$a>0$。}
\SubQ{3}{$\displaystyle\int_0^{\pi/2}\sin^6x\cos^4x\,dx$。}
\end{minipage}\par\vfill
\newpage
\noindent\begin{minipage}{\textwidth}\ProblemHeading{习题}{13.57}
利用 Euler 积分计算：
\SubQ{1}{$\displaystyle\int_0^1 x^{n-1}(1-x^m)^{q-1}\,dx$，$n,m,q>0$。}
\SubQ{2}{$\displaystyle\int_0^{+\infty}\frac{e^{-at}}{\sqrt{\pi t}}\,dt$，$a>0$。}
\SubQ{3}{$\displaystyle\int_0^{\pi/2}\tan^\alpha x\,dx$，$|\alpha|<1$。}
\end{minipage}\par\vfill
\noindent\begin{minipage}{\textwidth}\ProblemHeading{习题}{13.58}
利用 Euler 积分计算：
\SubQ{1}{$\displaystyle\int_0^1\sqrt{\frac1x\ln\frac1x}\,dx$。}
\SubQ{2}{$\displaystyle\int_a^b\left(\frac{b-x}{x-a}\right)^p\,dx$，$a<b$，$0<p<1$。}
\SubQ{3}{$\displaystyle\lim_{n\to\infty}\int_1^2(x-1)^2\sqrt[n]{\frac{2-x}{x-1}}\,dx$。}
\end{minipage}\par\vfill
\noindent\begin{minipage}{\textwidth}\ProblemHeading{习题}{13.59}
利用 Euler 积分计算：
\SubQ{1}{$\displaystyle\lim_{n\to\infty}\int_0^{+\infty}\frac{dx}{1+x^n}$。}
\end{minipage}\par\vfill
\noindent\begin{minipage}{\textwidth}\ProblemHeading{习题}{13.60}
计算极限 $\displaystyle\lim_{\alpha\to+\infty}\sqrt\alpha\int_0^1x^{3/2}(1-x^5)^\alpha\,dx$。
\end{minipage}\par\vfill
\newpage
\noindent\begin{minipage}{\textwidth}\ProblemHeading{习题}{13.61}
设 $a>0$，$n>0$。求曲线 $x^n+y^n=a^n$ 与两坐标轴围成的图形在第一象限的面积。\textbf{（校订：补明参数 $n>0$。）}
\end{minipage}\par\vfill
\noindent\begin{minipage}{\textwidth}\ProblemHeading{习题}{13.62}
设 $0<\alpha<1$。完成下列证明：
\SubQ{1}{对 $x>0$，证明 $x^\alpha-\alpha x+\alpha-1\le0$。}
\SubQ{2}{对 $a,b>0$，证明 $a^\alpha b^{1-\alpha}\le\alpha a+(1-\alpha)b$。}
\SubQ{3}{若 $x_i,y_i\ge0$（$1\le i\le n$），证明 Hölder 不等式 $\displaystyle\sum_{i=1}^nx_i y_i\le\left(\sum_{i=1}^nx_i^{1/\alpha}\right)^\alpha\left(\sum_{i=1}^ny_i^{1/(1-\alpha)}\right)^{1-\alpha}$。}
\end{minipage}\par\vfill
\newpage
\noindent\begin{minipage}{\textwidth}\ProblemHeading{习题}{13.63}
设 $0<\alpha<1$。完成下列证明：
\SubQ{1}{设 $f,g\ge0$ 且在 $[a,b]$ 上连续。利用有限和的 Hölder 不等式及 Riemann 和取极限，证明 $\displaystyle\int_a^bfg\,dx\le\left(\int_a^bf^{1/\alpha}\,dx\right)^\alpha\left(\int_a^bg^{1/(1-\alpha)}\,dx\right)^{1-\alpha}$。}
\SubQ{2}{设 $f,g\ge0$ 且在 $[a,+\infty)$ 上连续，所涉及反常积分均收敛。证明 $\displaystyle\int_a^{+\infty}fg\,dx\le\left(\int_a^{+\infty}f^{1/\alpha}\,dx\right)^\alpha\left(\int_a^{+\infty}g^{1/(1-\alpha)}\,dx\right)^{1-\alpha}$。}
\SubQ{3}{证明 $\ln\Gamma(x)$ 在 $(0,+\infty)$ 上为凸函数，其中 $\displaystyle\Gamma(x)=\int_0^{+\infty}t^{x-1}e^{-t}\,dt$。}
\end{minipage}\par\vfill
\newpage
\noindent\begin{minipage}{\textwidth}\ProblemHeading{习题}{13.64}
设 $a\ge0$，函数 $f\ge0$ 并在 $[a,+\infty)$ 的任意有限区间上可积。数列 $\{a_n\}$ 单调递增，$a_0=a$，$a_n\to+\infty$。证明 $\displaystyle\int_a^{+\infty}f(x)\,dx=l$ 当且仅当 $\displaystyle\sum_{n=1}^\infty\int_{a_{n-1}}^{a_n}f(x)\,dx=l$。
\end{minipage}\par\vfill
\noindent\begin{minipage}{\textwidth}\ProblemHeading{习题}{13.65}
证明 $\displaystyle\int_0^{+\infty}\frac{x}{1+x^6\sin^2x}\,dx$ 收敛，但被积函数 $\displaystyle f(x)=\frac{x}{1+x^6\sin^2x}$ 在 $[0,+\infty)$ 上非负、连续、无界，因而当 $x\to+\infty$ 时并不趋于 $0$。
\end{minipage}\par\vfill
\noindent\begin{minipage}{\textwidth}\ProblemHeading{习题}{13.66}
设 $\varphi$ 有二阶导数，$\psi$ 有一阶导数，$a>0$。证明\[u(x,t)=\frac12[\varphi(x-at)+\varphi(x+at)]+\frac1{2a}\int_{x-at}^{x+at}\psi(s)\,ds\]满足弦振动方程 $\displaystyle\frac{\partial^2u}{\partial t^2}=a^2\frac{\partial^2u}{\partial x^2}$。
\end{minipage}\par\vfill
\noindent\begin{minipage}{\textwidth}\ProblemHeading{习题}{13.67}
设 $n$ 为非负整数。证明 $n$ 阶 Bessel 函数 $\displaystyle J_n(x)=\frac1\pi\int_0^\pi\cos(n\varphi-x\sin\varphi)\,d\varphi$ 满足方程 $x^2J_n''(x)+xJ_n'(x)+(x^2-n^2)J_n(x)=0$。
\end{minipage}\par\vfill
\newpage
\noindent\begin{minipage}{\textwidth}\ProblemHeading{习题}{13.68}
证明：对任意实数 $u$，$\displaystyle\frac1{2\pi}\int_0^{2\pi}e^{u\cos x}\cos(u\sin x)\,dx=1$。
\end{minipage}\par\vfill
\noindent\begin{minipage}{\textwidth}\ProblemHeading{习题}{13.69}
证明积分 $\displaystyle\int_0^{+\infty}\frac{\sin3x}{x+u}e^{-ux}\,dx$ 关于 $u\in[0,+\infty)$ 一致收敛。
\end{minipage}\par\vfill
\noindent\begin{minipage}{\textwidth}\ProblemHeading{习题}{13.70}
设 $a>0$。证明积分 $\displaystyle\int_0^{+\infty}\frac{x\cos(ux)}{x^2+a^2}\,dx$ 关于 $u\in[\delta,+\infty)$（$\delta>0$）一致收敛，但关于 $u\in(0,+\infty)$ 不一致收敛。
\end{minipage}\par\vfill
\noindent\begin{minipage}{\textwidth}\ProblemHeading{习题}{13.71}
设 $a>0$。证明积分 $\displaystyle\int_0^{+\infty}\frac{x\sin(ux)}{x^2+a^2}\,dx$ 关于 $u\in(0,+\infty)$ 不一致收敛。
\end{minipage}\par\vfill
\newpage
\noindent\begin{minipage}{\textwidth}\ProblemHeading{习题}{13.72}
设 $\displaystyle\int_{-\infty}^{+\infty}|f(x)|\,dx$ 收敛。证明函数 $\displaystyle\varphi(u)=\int_{-\infty}^{+\infty}f(x)\cos(ux)\,dx$ 在 $\mathbb R$ 上一致连续。
\end{minipage}\par\vfill
\noindent\begin{minipage}{\textwidth}\ProblemHeading{习题}{13.73}
证明函数 $\displaystyle f(x)=\int_0^{+\infty}\frac{e^{-t}}{|\sin t|^x}\,dt$ 在 $[0,1)$ 上连续。
\end{minipage}\par\vfill
\noindent\begin{minipage}{\textwidth}\ProblemHeading{习题}{13.74}
证明 $\displaystyle\int_0^1\ln\Gamma(x)\,dx=\ln\sqrt{2\pi}$，其中 $\displaystyle\Gamma(x)=\int_0^{+\infty}t^{x-1}e^{-t}\,dt$（$x>0$）。
\end{minipage}\par\vfill
\noindent\begin{minipage}{\textwidth}\ProblemHeading{习题}{13.75}
证明 $\displaystyle\int_0^1\sin(\pi x)\ln\Gamma(x)\,dx=\frac1\pi\left(\ln\frac\pi2+1\right)$，其中 $\Gamma$ 为 Gamma 函数。
\end{minipage}\par\vfill
\newpage
\noindent\begin{minipage}{\textwidth}\ProblemHeading{习题}{13.76}
证明 $\displaystyle\int_0^\pi\frac{dx}{\sqrt{3-\cos x}}=\frac1{4\sqrt\pi}\Gamma^2\!\left(\frac14\right)$，其中 $\Gamma$ 为 Gamma 函数。
\end{minipage}\par\vfill
\noindent\begin{minipage}{\textwidth}\ProblemHeading{习题}{13.77}
设 $\varphi$ 是 $(0,+\infty)$ 上正、严格递减的连续函数，$\displaystyle\lim_{t\to0^+}\varphi(t)=+\infty$，$\displaystyle\int_0^{+\infty}\varphi(t)\,dt=1$。令 $\psi$ 为 $\varphi$ 的反函数。证明存在 $p\in(0,1)$，使 $\displaystyle\int_0^p\varphi(t)\,dt+\int_0^p\psi(t)\,dt=1+p^2$。
\end{minipage}\par\vfill
\noindent\begin{minipage}{\textwidth}\ProblemHeading{习题}{13.78}
设 $f$ 在 $[0,+\infty)$ 上连续，且 $\displaystyle\int_0^{+\infty}f^2(x)\,dx$ 收敛。证明函数 $\displaystyle g(x)=f(x)-2e^{-x}\int_0^xe^tf(t)\,dt$ 满足 $\displaystyle\int_0^{+\infty}g^2(x)\,dx=\int_0^{+\infty}f^2(x)\,dx$。
\end{minipage}\par\vfill
